\DocumentMetadata{
  pdfversion=2.0,
  pdfstandard=ua-2,
  lang=en-US,
  tagging=on,
  tagging-setup={math/setup=mathml-SE,tabsorder=structure}
}
\documentclass[11pt,letterpaper]{article}
\usepackage[margin=0.68in,includefoot]{geometry}
\usepackage{newcomputermodern}
\usepackage{amsmath,amscd,mathtools}
\usepackage{tikz,adjustbox}
\providecommand{\square}{\mdlgwhtsquare}
\usepackage[hidelinks]{hyperref}
\hypersetup{
  pdftitle={Algebra First-Year Exam, May 2017, Part 2},
  pdfauthor={University of Florida Department of Mathematics},
  pdfsubject={Graduate mathematics examination},
  pdfdisplaydoctitle=true,
  bookmarksopen=true
}
\setcounter{secnumdepth}{0}
\setlength{\parindent}{0pt}
\setlength{\parskip}{0.28em}
\setlength{\emergencystretch}{3em}
\setlength{\leftmargini}{2.15em}
\setlength{\leftmarginii}{2.65em}
\setlength{\leftmarginiii}{3em}
\pagestyle{empty}
\raggedbottom
\allowdisplaybreaks
\NewTaggingSocketPlug{block/list/label}{examproblem}{%
  \tagstructbegin{tag=Lbl}\tagstructend%
  \tagstructbegin{tag=\LBody}%
  \tagstructbegin{tag=\ProblemHeadingTag,title-o={\ProblemHeadingText}}%
  \tagmcbegin{tag=\ProblemHeadingTag,actualtext={\ProblemHeadingText}}%
  #2%
  \tagmcend%
  \tagstructend%
}
\newenvironment{examproblems}[2]{%
  \def\ProblemHeadingTag{#2}%
  \AssignTaggingSocketPlug{block/list/label}{examproblem}%
  \begin{enumerate}%
  \setcounter{enumi}{#1}%
  \renewcommand{\labelenumi}{\textbf{\theenumi.}}%
  \setlength{\topsep}{0.35em}%
  \setlength{\partopsep}{0pt}%
  \setlength{\itemsep}{0.48em}%
  \setlength{\parsep}{0.18em}%
}{\end{enumerate}}
\newenvironment{examsubparts}{%
  \AssignTaggingSocketPlug{block/list/label}{default}%
  \begin{enumerate}%
  \setlength{\topsep}{0.18em}%
  \setlength{\partopsep}{0pt}%
  \setlength{\itemsep}{0.16em}%
  \setlength{\parsep}{0.1em}%
}{\end{enumerate}}
\newenvironment{examinstructions}{%
  \par\begingroup\itshape
}{%
  \par\endgroup\vspace{0.18em}
}
\makeatletter
\renewcommand\subsection{\@startsection{subsection}{2}{0pt}%
  {-1.25ex plus -.25ex minus -.1ex}%
  {0.55ex plus .1ex}%
  {\normalfont\large\bfseries}}
\renewcommand{\@maketitle}{%
  \newpage\null\vskip -1em%
  {\footnotesize\bfseries
    \href{https://www.ufl.edu/}{University of Florida}, \href{https://math.ufl.edu/}{Department of Mathematics}\par}%
  \vskip -0.5em%
  \tagstructbegin{tag=H1,title={Algebra First-Year Exam, May 2017, Part 2}}%
  \tagpdfparaOff%
  \tagmcbegin{tag=H1}%
    {\LARGE\bfseries\href{https://gma.math.ufl.edu/exams/algebra-first-year/}{Algebra First-Year Exam}, May 2017, Part 2\par}%
  \tagmcend%
  \tagpdfparaOn%
  \tagstructend%
  \vskip 0.25em%
  \tagmcbegin{artifact}%
    \hrule height 1pt%
  \tagmcend%
  \vskip 1em%
}
\makeatother
\title{Algebra First-Year Exam, May 2017, Part 2}
\author{}
\date{}
\begin{document}
\maketitle
\thispagestyle{empty}
\begin{examinstructions}
Answer four problems. If you turn in more than four, only the first four will be graded. Within reason, you may use theorems as long as you state them clearly.
\end{examinstructions}
\begin{examproblems}{0}{H2}
\def\ProblemHeadingText{Problem 1.}
\item
Let \(\omega=\frac{-1+\sqrt{-3}}{2}\). Show that the ring \(\mathbb{Z}[\omega]\) of all complex numbers of the form \(a+b\omega\) with \(a,b\in\mathbb{Z}\) is a Euclidean domain.
\def\ProblemHeadingText{Problem 2.}
\item
Let~\(F\) be a field.
\begin{examsubparts}
\item[\textnormal{(i)}]
Show that if~\(F\) has characteristic zero and \(n\geq 1\) is an integer, the ring \(F[X,Y]/(X^n+Y^n-1)\) is an integral domain. (Hint: first show that \(Y-1\) and \(Y^{n-1}+Y^{n-2}+\cdots+Y+1\) are relatively prime in \(F[Y]\).)
\item[\textnormal{(ii)}]
Find an example of a field of positive characteristic and \(n\geq 1\) for which the conclusion of (i) is false.
\end{examsubparts}
\def\ProblemHeadingText{Problem 3.}
\item
Let~\(R\) be an integral domain. Recall that an~\(R\)-module~\(M\) has rank~\(n\) if~\(n\) is the largest integer such that~\(M\) has~\(n\) \(R\)-linearly independent elements. Show that an~\(R\)-module~\(M\) has rank~\(n\) if and only if~\(M\) has a submodule \(N\subseteq M\) such that~\(N\) is free of rank~\(n\), and \(M/N\) is torsion.
\def\ProblemHeadingText{Problem 4.}
\item
Let~\(F\) be a field, \(M\) an extension of~\(F\) and \(L,K\) two finite extensions of~\(F\) contained in~\(M\). Show that if \([L:F]\) and \([K:F]\) are relatively prime, \(L\cap K=F\).
\def\ProblemHeadingText{Problem 5.}
\item
Show that the group \(GL_3(\mathbb{F}_2)\) has two distinct conjugacy classes of elements of order~\(7\), and find representatives for each class.
\def\ProblemHeadingText{Problem 6.}
\item
Let~\(F\) be a field, \(V\) an~\(F\)-vector space of finite dimension and \(T:V\to V\) a linear transformation. Show that if the minimal polynomial of~\(T\) is irreducible of degree~\(d\) then the dimension of~\(V\) is a multiple of~\(d\).
\end{examproblems}
\end{document}
